\section{13. Risks and critical reflection}

\textbf{The headline result is that the clever part buys little.} The proposal
set out to show that predicting occupancy pays off. It does not, here: reactive
control on \COtwo{} is within four room-minutes per week of perfect foresight,
and our predictor closes that gap for 3.8~\% more energy. We think the finding
is real for this model and would not carry to every building. The dampers reach
maximum flow in minutes and the maximum is 1.5 times design flow, so a
reactive controller is never far behind. A building with slower or smaller
ventilation would show a larger gap. We did not vary those constants, which is
the experiment we would run next.

\textbf{The predictor cannot see what matters most.} The rooms that fill
fastest are the lecture rooms, and they are the ones with no pattern: the
forecast foresaw none of their arrivals. The fix is an input, not a better
model: the day's bookings are available from occupancysim and could be read as
a booking system would provide them. We parked that for time.

\textbf{We measured the forecast with the wrong yardstick.} FR-4 asked for a
lower mean error than ``same as now'', which a deliberately high forecast
cannot give. We report it as not met and say what should have been asked:
arrivals foreseen at a stated false-alarm rate.

\textbf{What changed from the proposal, and what it cost.} The proposal's main
claim was occupancy without an occupancy sensor. We gave each room a people
counter instead, and the system no longer shows that \COtwo{} alone is enough.
The mass balance we meant to invert is used only as two bounds, to catch a
counter that is stuck. A full estimator would also give the degraded mode a
head-count when the counter fails, which today it replaces with the forecast.

\textbf{Mistakes that the tests did not catch.} Three faults got past the
unit tests: the predictor that learned sensor noise (its tests used clean
data), the runs of one day that were not identical (we had written ``same
seed'' under fairness without checking), and FR-7, which promised a fallback
for a stuck sensor that nothing detected. The first two were found by looking
at evaluation results that seemed implausible, the third by re-reading what we
had promised while writing the alerts. None was found by a failing test. Each
now has a regression test or a recorded run.

\begin{longtable}[]{@{}L{0.30\linewidth}L{0.28\linewidth}L{0.34\linewidth}@{}}
\toprule
\textbf{Risk or assumption} & \textbf{Impact} & \textbf{Mitigation, or decision} \\
\midrule
\endhead
The room physics is our own model, with assumed constants & Energy and comfort figures say as much about the model as about the controllers & Reported as comparisons between policies on the same model, never as absolute savings \\
A drifting sensor is not detected & A \COtwo{} sensor drifting up ends in the override: safe, wasteful. One drifting down would be believed & Not solved. Cross-checking neighbouring rooms, or the counter against \COtwo{}, is the next step \\
A people counter stuck on a small number is not found & Seen live: stuck at 4 in a 25-person room for three hours & Accepted: it costs the airflow for four people. Pinned by a unit test so the limit is known \\
A stuck counter is found only after the air disagrees & 15--20 minutes in a lecture room, about 75 in an office & Accepted; the cost is airflow, not air quality \\
The stuck-sensor limits come from our own noise model & A real sensor quieter than its resolution repeats itself when healthy and would be blamed & The limits would have to be set from the real sensor's noise at commissioning \\
The truth API is closed to the controller by convention only & The controller could call physics; it has no code that does & The truth \emph{topic} is enforced by the broker. A separate network for the truth layer would close the rest \\
Devices do not authenticate to the registry & Any process on the network can pose as a device & A real installation would give each device a secret at commissioning \\
The broker is a single point for observations & Decisions continue through the BuildSim fallback, but storage records nothing while it is down & Accepted and reported: a 20-minute gap in T-24 \\
The live count reacts to one false reading & About 5~\% of the predictive policy's decisions in an empty room ventilate for ``one person'' & Requiring two readings would remove it, at a minute's delay for a real first arrival. Not done \\
Command latency was never measured & NFR-4 is unverified & Records carry model time only; a real-time stamp is needed first \\
One seed and one week & The ranking of the policies is untested on other weeks & A second seed is five more hours of runs \\
250 containers on one laptop & A stock laptop stops at about 40 rooms; with three host settings raised, 50 run. More was not tried & One process per device is right for a real installation and expensive to simulate \\
The heating's warm-up rate is assumed & Rooms are warm about an hour early & Learning the rate per room from earlier mornings \\
\bottomrule
\end{longtable}

\textbf{How the design was reviewed.} We used an AI assistant throughout, and
kept a log of where its advice was wrong and how we caught it
(\code{docs/review-log.md}). The pattern is worth stating: its first designs
were plausible and broke rules it had not been shown. It put the sensor noise
inside physics, shared files between processes and showed true \COtwo{} in the
viewer, all three against the lectures; it treated occupancysim as something
we could not change; and it proposed the stuck-counter rule that blames every
healthy counter. Each was caught by checking the design against the course
material or by asking what a healthy system looks like, before or after code
was written. The parts of this design we are most sure of are the ones that
were wrong once.

\todo{Decide whether the paragraph above belongs in the report, and check
what Canvas says about declaring AI use. This draft was itself written with
the assistant, from your documents.}

\textbf{With more time.} In order: vary the damper and room constants to see
when anticipation starts to pay; give the predictor the timetable; measure
command latency; a second seed; drift detection; and the \COtwo{}-based
occupancy estimate the proposal started from.
